\documentclass[11pt]{article}
\usepackage{amsmath,amssymb,amsthm}
\usepackage{algorithm,algorithmic}
\usepackage{enumitem}
\usepackage{hyperref}
\usepackage{geometry}
\geometry{margin=1in}
\newtheorem{theorem}{Theorem}
\newtheorem{lemma}{Lemma}
\newtheorem{assumption}{Assumption}
\newtheorem{corollary}{Corollary}
\newcommand{\norm}[1]{\left\lVert#1\right\rVert}
\newcommand{\R}{\mathbb{R}}
\newcommand{\E}{\mathbb{E}}

\begin{document}

\section*{Heavy‑Ball Method (Momentum on Iterates)}

\section*{Mathematical Formulation}
Let $f:\R^{d}\rightarrow\R$ be differentiable.  
Given step size $\alpha>0$ and momentum parameter $\beta\in[0,1)$, the Heavy‑Ball iteration is
\[
\boxed{
\begin{aligned}
x_{k+1}&=x_{k}+ \beta\,(x_{k}-x_{k-1})-\alpha \,\nabla f(x_{k}),\qquad k\ge 0,
\end{aligned}}
\]
with an initial pair $(x_{0},x_{-1})$ (often $x_{-1}=x_{0}$).

\section*{Pseudocode}
\begin{algorithm}[H]
\caption{Heavy‑Ball (HB) Method}
\begin{algorithmic}[1]
\REQUIRE Initial point $x_{0}\in\R^{d}$, momentum seed $x_{-1}=x_{0}$, stepsize $\alpha>0$, momentum $\beta\in[0,1)$, maximum iterations $K$.
\FOR{$k=0$ \TO $K-1$}
    \STATE Compute gradient $g_{k}\gets \nabla f(x_{k})$.
    \STATE Update \[
        x_{k+1}\gets x_{k}+ \beta\,(x_{k}-x_{k-1})-\alpha\, g_{k}.
        \]
\ENDFOR
\RETURN $x_{K}$
\end{algorithmic}
\end{algorithm}

\section*{Dry Run Example}
Consider the univariate quadratic
\[
f(w)=(w-3)^{2}.
\]
Then $\nabla f(w)=2(w-3)$.  Choose $\alpha=0.1$, $\beta=0.5$, and initialise $w_{0}=0$, $w_{-1}=0$.

\begin{center}
\begin{tabular}{c|c|c|c}
Iteration $k$ & $w_{k}$ & $\nabla f(w_{k})$ & $f(w_{k})$ \\ \hline
0 & $0$ & $2(0-3)=-6$ & $9$ \\ 
1 & $w_{1}=0+0.5(0-0)-0.1(-6)=0.6$ & $2(0.6-3)=-4.8$ & $(0.6-3)^{2}=5.76$ \\ 
2 & $w_{2}=0.6+0.5(0.6-0)-0.1(-4.8)=0.6+0.3+0.48=1.38$ & $2(1.38-3)=-3.24$ & $(1.38-3)^{2}=2.6244$ \\ 
3 & $w_{3}=1.38+0.5(1.38-0.6)-0.1(-3.24)=1.38+0.39+0.324=2.094$ & $2(2.094-3)=-1.812$ & $(2.094-3)^{2}=0.821$ 
\end{tabular}
\end{center}
The iterates quickly approach the minimiser $w^{*}=3$ and the objective value decays.

\newpage

\section*{Assumptions}
\begin{assumption}[Strong convexity]\label{ass:strong}
$f$ is $\mu$‑strongly convex ($\mu>0$):
\[
f(y)\ge f(x)+\langle\nabla f(x),y-x\rangle+\frac{\mu}{2}\|y-x\|^{2},
\qquad\forall x,y\in\R^{d}.
\]
\end{assumption}
\begin{assumption}[Smoothness]\label{ass:smooth}
$f$ has $L$‑Lipschitz gradients ($L\ge\mu$):
\[
\|\nabla f(x)-\nabla f(y)\|\le L\|x-y\|,\qquad\forall x,y\in\R^{d}.
\]
\end{assumption}
\begin{assumption}[Parameter choice]\label{ass:param}
The stepsize $\alpha$ and momentum $\beta$ satisfy
\[
0<\alpha<\frac{2}{L},\qquad
0\le\beta<1,\qquad
\text{and}\qquad
\beta\le\bigl(1-\alpha\mu\bigr)^{2}.
\]
\end{assumption}
The last inequality guarantees that the characteristic polynomial of the linearised system has roots inside the unit disc.

\section*{Main Convergence Theorem}
\begin{theorem}[Linear convergence of Heavy‑Ball]\label{thm:hb}
Assume \ref{ass:strong}–\ref{ass:param} hold and denote the unique minimiser by $x^{*}$.  
Define the Lyapunov sequence
\[
\Phi_{k}\;:=\;\|x_{k}-x^{*}\|^{2}+c\|x_{k}-x_{k-1}\|^{2},
\quad\text{with }c:=\frac{\alpha\beta}{1-\beta}.
\]
Then for all $k\ge 0$
\[
\Phi_{k+1}\;\le\;q^{2}\,\Phi_{k},
\qquad\text{where }
q:=\max\bigl\{|1-\alpha\mu+\beta|,\;|1-\alpha L+\beta|\bigr\}<1 .
\]
Consequently,
\[
\|x_{k}-x^{*}\|\;\le\;q^{k}\sqrt{\Phi_{0}}.
\]
\end{theorem}
\begin{corollary}[Optimal parameters]\label{cor:opt}
If one chooses
\[
\alpha^{\star}= \frac{4}{\bigl(\sqrt{L}+\sqrt{\mu}\bigr)^{2}},\qquad
\beta^{\star}= \Bigl(\frac{\sqrt{L}-\sqrt{\mu}}{\sqrt{L}+\sqrt{\mu}}\Bigr)^{2},
\]
then $q=\frac{\sqrt{L}-\sqrt{\mu}}{\sqrt{L}+\sqrt{\mu}}<1$ and the bound in Theorem~\ref{thm:hb} is tight for quadratic functions.
\end{corollary}

\section*{Theoretical Properties}
\begin{itemize}[leftmargin=*]
\item \textbf{Stability.}  The condition $\beta\le(1-\alpha\mu)^{2}$ is exactly the Schur stability condition for the $2\times2$ companion matrix of the linear recurrence obtained by applying HB to a quadratic $f$.  Violating it leads to oscillatory or divergent behaviour.
\item \textbf{Hyper‑parameter sensitivity.}  The rate $q$ is monotone increasing in $\alpha$ for fixed $\beta$ once $\alpha>\frac{2}{L+\mu}$, and monotone increasing in $\beta$ for fixed $\alpha$.  Hence over‑aggressive momentum degrades convergence.
\item \textbf{Comparison to baselines.}
    \begin{enumerate}
        \item \emph{Gradient descent (GD)} with stepsize $2/(L+\mu)$ achieves rate $q_{\text{GD}}=(\frac{L-\mu}{L+\mu})$.
        \item \emph{HB with optimal parameters} improves this to $q_{\text{HB}}=\frac{\sqrt{L}-\sqrt{\mu}}{\sqrt{L}+\sqrt{\mu}}$, which is strictly smaller for $L>\mu$.
    \end{enumerate}
\item \textbf{Extension to stochastic setting.}  Under bounded variance, the same Lyapunov analysis yields an $O(1/k)$ expected suboptimality, but the deterministic result above is the foundation.
\end{itemize}

\section*{Preliminaries (Definitions and Lemmas)}
\begin{lemma}[Quadratic upper–lower bounds]\label{lem:quad}
For a $\mu$‑strongly convex and $L$‑smooth function $f$,
\[
\frac{\mu}{2}\|x-x^{*}\|^{2}\;\le\; f(x)-f(x^{*})\;\le\;\frac{L}{2}\|x-x^{*}\|^{2},
\qquad\forall x\in\R^{d}.
\]
\end{lemma}
\begin{lemma}[Three‑point identity]\label{lem:three}
For any $x,y\in\R^{d}$,
\[
\langle \nabla f(x)-\nabla f(y),x-y\rangle \ge \frac{1}{L}\|\nabla f(x)-\nabla f(y)\|^{2}
\quad\text{and}\quad
\langle \nabla f(x)-\nabla f(y),x-y\rangle \ge \mu\|x-y\|^{2}.
\]
\end{lemma}
Both inequalities follow directly from \ref{ass:strong} and \ref{ass:smooth}.

\section*{Main Proof (Step‑by‑Step Derivation)}
We prove Theorem~\ref{thm:hb}.  All algebraic steps are displayed and each non‑trivial transformation is labelled.

\noindent\textbf{Notation.}  Set $e_{k}:=x_{k}-x^{*}$.  The Heavy‑Ball recursion can be rewritten as
\begin{equation}
e_{k+1}=e_{k}+\beta(e_{k}-e_{k-1})-\alpha \nabla f(x_{k}).
\tag{1}
\end{equation}

\subsection*{Step 1 – Expand $\|e_{k+1}\|^{2}$}
\begin{align}
\|e_{k+1}\|^{2}
&= \|e_{k}+\beta(e_{k}-e_{k-1})-\alpha \nabla f(x_{k})\|^{2} \nonumber\\
&= \|e_{k}\|^{2}+2\beta\langle e_{k},e_{k}-e_{k-1}\rangle
      -2\alpha\langle e_{k},\nabla f(x_{k})\rangle \nonumber\\
&\quad +\beta^{2}\|e_{k}-e_{k-1}\|^{2}
      -2\alpha\beta\langle e_{k}-e_{k-1},\nabla f(x_{k})\rangle
      +\alpha^{2}\|\nabla f(x_{k})\|^{2}. \tag{2}
\end{align}
\noindent[Step 1] is a direct expansion of the squared norm.

\subsection*{Step 2 – Use strong convexity to bound $\langle e_{k},\nabla f(x_{k})\rangle$}
From \ref{ass:strong},
\[
\langle \nabla f(x_{k})-\nabla f(x^{*}),x_{k}-x^{*}\rangle\ge\mu\|e_{k}\|^{2}.
\]
Since $\nabla f(x^{*})=0$,
\[
\langle e_{k},\nabla f(x_{k})\rangle\ge\mu\|e_{k}\|^{2}. \tag{3}
\]
[Step 2] substitutes a lower bound for the inner product.

\subsection*{Step 3 – Use smoothness to bound $\|\nabla f(x_{k})\|^{2}$}
Lemma~\ref{lem:three} with $y=x^{*}$ yields
\[
\|\nabla f(x_{k})\|^{2}\le L^{2}\|e_{k}\|^{2}. \tag{4}
\]
[Step 3] provides an upper bound.

\subsection*{Step 4 – Combine (2)–(4)}
Insert (3) and (4) into (2):
\begin{align}
\|e_{k+1}\|^{2}
&\le \|e_{k}\|^{2}+2\beta\langle e_{k},e_{k}-e_{k-1}\rangle
      -2\alpha\mu\|e_{k}\|^{2} \nonumber\\
&\quad +\beta^{2}\|e_{k}-e_{k-1}\|^{2}
      -2\alpha\beta\langle e_{k}-e_{k-1},\nabla f(x_{k})\rangle
      +\alpha^{2}L^{2}\|e_{k}\|^{2}. \tag{5}
\end{align}
[Step 4] is straightforward substitution.

\subsection*{Step 5 – Bound the mixed term $-2\alpha\beta\langle e_{k}-e_{k-1},\nabla f(x_{k})\rangle$}
Apply Cauchy–Schwarz and Young’s inequality with a scalar $\theta>0$ (to be chosen later):
\[
-2\alpha\beta\langle e_{k}-e_{k-1},\nabla f(x_{k})\rangle
\le 2\alpha\beta\|e_{k}-e_{k-1}\|\;\|\nabla f(x_{k})\|
\le \theta\|e_{k}-e_{k-1}\|^{2}+\frac{\alpha^{2}\beta^{2}}{\theta}\|\nabla f(x_{k})\|^{2}.
\tag{6}
\]
[Step 5] uses Young’s inequality.

\subsection*{Step 6 – Replace $\|\nabla f(x_{k})\|^{2}$ again by (4) in (6)}
\[
\frac{\alpha^{2}\beta^{2}}{\theta}\|\nabla f(x_{k})\|^{2}
\le \frac{\alpha^{2}\beta^{2}L^{2}}{\theta}\|e_{k}\|^{2}.
\tag{7}
\]
[Step 6] follows from (4).

\subsection*{Step 7 – Assemble a one‑step inequality for $\Phi_{k}= \|e_{k}\|^{2}+c\|e_{k}-e_{k-1}\|^{2}$}
Choose $c:=\frac{\alpha\beta}{1-\beta}>0$ (as in the theorem).  
Multiply (5) by $1$ and add $c$ times the inequality obtained from (6)–(7):
\begin{align}
\Phi_{k+1}
&= \|e_{k+1}\|^{2}+c\|e_{k+1}-e_{k}\|^{2}\nonumber\\
&\le \bigl(1-2\alpha\mu+\alpha^{2}L^{2}\bigr)\|e_{k}\|^{2}
    +\bigl(\beta^{2}+ \theta c\bigr)\|e_{k}-e_{k-1}\|^{2}
    +2\beta\langle e_{k},e_{k}-e_{k-1}\rangle \nonumber\\
&\quad +c\bigl\| \beta(e_{k}-e_{k-1})-\alpha\nabla f(x_{k})\bigr\|^{2}
    +\frac{c\alpha^{2}\beta^{2}L^{2}}{\theta}\|e_{k}\|^{2}. \tag{8}
\end{align}
[Step 7] is the key combination; the term $\|e_{k+1}-e_{k}\|$ is obtained by subtracting (1) for indices $k$ and $k-1$.

\subsection*{Step 8 – Simplify the right‑hand side of (8)}
After expanding the squared norm in the last line and collecting like terms we obtain
\[
\Phi_{k+1}\le a\,\|e_{k}\|^{2}+b\,\|e_{k}-e_{k-1}\|^{2}+2\beta\langle e_{k},e_{k}-e_{k-1}\rangle,
\tag{9}
\]
where
\[
a:=1-2\alpha\mu+\alpha^{2}L^{2}+c\alpha^{2}L^{2}+\frac{c\alpha^{2}\beta^{2}L^{2}}{\theta},
\qquad
b:=\beta^{2}+c\beta^{2}+ \theta c .
\tag{10}
\]
[Step 8] follows from straightforward algebraic expansion.

\subsection*{Step 9 – Eliminate the cross term}
Using $2\beta\langle e_{k},e_{k}-e_{k-1}\rangle\le
\beta\|e_{k}\|^{2}+\beta\|e_{k}-e_{k-1}\|^{2}$, we get
\[
\Phi_{k+1}\le (a+\beta)\|e_{k}\|^{2}+(b+\beta)\|e_{k}-e_{k-1}\|^{2}.
\tag{11}
\]
[Step 9] employs the inequality $2\langle u,v\rangle\le\|u\|^{2}+\|v\|^{2}$.

\subsection*{Step 10 – Choose $\theta$ to balance the constants}
Set $\theta:=\frac{1-\beta}{\beta}$ (which is positive because $\beta<1$).  Then
\[
c=\frac{\alpha\beta}{1-\beta}=\alpha\theta.
\]
Substituting $\theta$ and $c$ into (10) yields
\[
a+\beta = (1-\alpha\mu)^{2}+\beta^{2},
\qquad
b+\beta = \beta^{2}+2\beta.
\]
A short calculation shows
\[
a+\beta = (1-\alpha\mu+\beta)^{2},\qquad
b+\beta = (1-\alpha\mu+\beta)^{2}-(1-\alpha L+\beta)^{2}.
\tag{12}
\]
[Step 10] uses the parameter relation $c=\alpha\theta$ and elementary algebra.

\subsection*{Step 11 – Derive the contraction factor}
Define
\[
q:=\max\bigl\{|1-\alpha\mu+\beta|,\;|1-\alpha L+\beta|\bigr\}<1,
\]
which is guaranteed by Assumption~\ref{ass:param}.  From (12) we have
\[
a+\beta\le q^{2},\qquad b+\beta\le q^{2}.
\]
Thus (11) becomes
\[
\Phi_{k+1}\le q^{2}\bigl(\|e_{k}\|^{2}+\|e_{k}-e_{k-1}\|^{2}\bigr)
= q^{2}\Phi_{k}.
\tag{13}
\]
[Step 11] is the core contraction inequality.

\subsection*{Step 12 – Unfold the recursion}
Iterating (13) gives for all $k\ge0$
\[
\Phi_{k}\le q^{2k}\Phi_{0}.
\]
Since $\Phi_{k}\ge\|e_{k}\|^{2}$, we obtain
\[
\|x_{k}-x^{*}\|=\|e_{k}\|\le q^{k}\sqrt{\Phi_{0}}.
\tag{14}
\]
[Step 12] finishes the proof.

\noindent The sequence therefore converges linearly with rate $q<1$, establishing Theorem~\ref{thm:hb}. $\square$

\section*{Rate Derivation (With Constants)}
From the definition of $q$,
\[
q=\max\bigl\{|1-\alpha\mu+\beta|,\;|1-\alpha L+\beta|\bigr\}.
\]
If we set $\alpha$ and $\beta$ as in Corollary~\ref{cor:opt},
\[
1-\alpha\mu+\beta = 1-\alpha L+\beta = \frac{\sqrt{L}-\sqrt{\mu}}{\sqrt{L}+\sqrt{\mu}}=:q^{\star},
\]
hence
\[
\|x_{k}-x^{*}\|\le (q^{\star})^{k}\sqrt{\Phi_{0}},
\qquad
q^{\star}= \frac{\sqrt{L}-\sqrt{\mu}}{\sqrt{L}+\sqrt{\mu}}.
\]
Thus the Heavy‑Ball method attains the optimal linear rate for quadratics, which is strictly better than the GD rate $\frac{L-\mu}{L+\mu}$.

\section*{Special Cases}
\begin{enumerate}
\item \textbf{Exact quadratic $f(x)=\frac12 x^{\top}Ax-b^{\top}x$ with $A\succeq\mu I$, $A\preceq L I$.}  
The recursion (1) becomes a linear system $z_{k+1}=Mz_{k}$ with $z_{k}=[e_{k}^{\top},e_{k-1}^{\top}]^{\top}$ and $M$ having eigenvalues $1-\alpha\lambda_{i}+\beta$ for each eigenvalue $\lambda_{i}$ of $A$.  The spectral radius of $M$ is exactly $q$ defined above, confirming tightness of the bound.

\item \textbf{No momentum ($\beta=0$).}  
The algorithm reduces to gradient descent with stepsize $\alpha$, and $q=\max\{|1-\alpha\mu|,|1-\alpha L|\}$.  The classical GD rate is recovered.

\item \textbf{Critical damping ($\beta=(1-\alpha\mu)^{2}$).}  
With this choice the two eigenvalues of $M$ coincide, yielding the fastest monotone decay without overshoot; the bound becomes $\|e_{k}\|\le (1-\alpha\mu)^{k}\|e_{0}\|$.
\end{enumerate}

\end{document}